\section{TDD: rojo-verde-refactorización}

\subsection{Disciplina de tres etapas}

\begin{importantbox}{Rojo-verde-refactorización}
\begin{enumerate}
    \item \textbf{Rojo (Red)}: primero se escribe una prueba que \textbf{fallará}; como la funcionalidad todavía no está implementada, la prueba necesariamente marca error.
    \item \textbf{Verde (Green)}: se escribe la \textbf{cantidad mínima} de código para que la prueba pase.
    \item \textbf{Refactorización (Refactor)}: se optimiza la estructura del código manteniéndolo en verde.
\end{enumerate}
\end{importantbox}

\subsection{¿Por qué es indispensable "ver con los propios ojos que la prueba falla"?}

\begin{warningbox}{Una prueba escrita después no demuestra nada}
\begin{itemize}
    \item Después de escribir la prueba hay que ejecutarla y ver el \textbf{error esperado}.
    \item Si la prueba pasa de inmediato, significa que lo que se está probando está mal, o que se está probando una conducta que ya existía.
    \item Escribir primero el código y agregar la prueba después casi siempre solo prueba "la conducta actual del código", no "la conducta que el código debería tener"; de ahí viene la \textbf{deuda técnica}.
\end{itemize}
\end{warningbox}

\subsection{La verificación necesita evidencia}

Antes de declarar que algo está terminado o que una prueba pasó, hay que \textbf{ejecutar de verdad el comando o la herramienta} y revisar la salida, usando abundantes aserciones (assert) con una expectativa clara sobre el resultado.

\subsection{Resumen de la sección}

La disciplina central de TDD es "primero rojo, luego verde, después refactorización": obliga al Agent a usar en cada paso resultados reales de pruebas como evidencia, y elimina el "fingir que algo está terminado".

